\documentclass[]{article}

\usepackage{amsmath}
\usepackage{multirow}
\usepackage{natbib}
\usepackage{graphicx} 


\begin{document}

We present the results of our numerical investigation into tracer transport in quenched 3D vortex filament fields. We first provide a quantitative summary of the key transport parameters, followed by a detailed analysis of the mean squared displacement, velocity statistics, temporal correlations, and the connection between transport dynamics and the local flow topology.

\subsubsection{Summary of transport parameters}

Tables \ref{tab:params_vortex}, \ref{tab:params_levy}, and \ref{tab:params_ow} consolidate the primary quantitative findings from our simulations. Table \ref{tab:params_vortex} lists the anomalous diffusion exponents $\alpha(N)$, velocity correlation times $\tau_c$, and persistence lengths $L_p$ for each filament density $N$. For comparison, Table \ref{tab:params_levy} provides the measured diffusion exponents for canonical isotropic L\'evy walks. Table \ref{tab:params_ow} summarizes the correlation between tracer speed and local flow topology for the highest density case, $N=40$.

\begin{table}[h!]
\centering
\caption{Summary of anomalous transport parameters for the 3D vortex filament system. Columns include the number of filaments ($N$), the fitted anomalous diffusion exponent ($\alpha(N)$) with its standard error (SE), the velocity autocorrelation zero-crossing time ($\tau_c$), the persistence length ($L_p$), and the mean inter-filament spacing ($\ell_{\text{inter}}$). Dashes indicate cases where the VACF did not cross zero within the observation time.}
\label{tab:params_vortex}
\begin{tabular}{cccccc}
\hline
$N$ & $\alpha(N)$ & $\alpha$ SE & $\tau_c$ (s) & $L_p$ (m) & $\ell_{\text{inter}}$ (m) \\
\hline
5  & 1.999 & 0.000 & —     & —     & 11.70 \\
10 & 1.964 & 0.002 & 10.17 & 0.313 & 9.28  \\
20 & 1.969 & 0.001 & —     & —     & 7.37  \\
40 & 1.835 & 0.003 & 9.555 & 0.578 & 5.85  \\
\hline
\end{tabular}
\end{table}

\begin{table}[h!]
\centering
\caption{Summary of transport parameters for the 3D isotropic L\'evy walk ground truth. Columns show the flight time distribution exponent ($\beta$), the theoretical diffusion exponent ($\alpha_{\text{theory}}$), and the exponent measured from our finite-time simulations ($\alpha_{\text{measured}}$).}
\label{tab:params_levy}
\begin{tabular}{cccc}
\hline
$\beta$ & $\alpha_{\text{theory}}$ & $\alpha_{\text{measured}}$ & $\tau_c$ (s) \\
\hline
1.2 & 1.8 & 1.655 & — \\
1.5 & 1.5 & 1.049 & — \\
1.8 & 1.2 & 1.069 & — \\
2.5 & 1.0 & 1.012 & — \\
\hline
\end{tabular}
\end{table}

\begin{table}[h!]
\centering
\caption{Okubo-Weiss and spatial correlation analysis for the $N=40$ configuration.}
\label{tab:params_ow}
\begin{tabular}{lc}
\hline
Quantity & Value \\
\hline
OW–speed Pearson r & −0.508 \\
Rolling speed variance–speed Pearson r & +0.425 \\
\hline
\end{tabular}
\end{table}

\subsection{Anomalous diffusion and MSD scaling}

The ensemble-averaged time-averaged mean squared displacement (TAMSD) curves for all four vortex filament configurations exhibit clear power-law growth, $\langle \delta^2(\Delta t) \rangle_t \sim (\Delta t)^\alpha$, over intermediate lag times. The fitted anomalous diffusion exponents $\alpha(N)$, reported in Table \ref{tab:params_vortex}, reveal a strongly superdiffusive regime that systematically evolves with filament density.

For the sparsest configuration ($N=5$), the transport is nearly ballistic, with a fitted exponent $\alpha = 1.999 \pm 0.000$. As the density increases to $N=10$ and $N=20$, the exponents remain high ($\alpha \approx 1.96-1.97$), indicating that transport is still dominated by long, persistent flights. Only at the highest density ($N=40$) does the exponent decrease significantly to $\alpha = 1.835 \pm 0.003$. This trend suggests a slow convergence towards the asymptotic regime predicted by Holtsmark theory, where $\alpha = 1.5$. At low densities, a tracer is primarily influenced by the nearest filament, leading to prolonged, nearly straight-line motion. As $N$ increases, the superposition of fields from multiple filaments introduces the velocity randomization necessary for anomalous, rather than ballistic, diffusion. However, even at $N=40$, the system remains far from the theoretical limit, highlighting the persistence of velocity correlations.

This behavior contrasts with the L\'evy walk simulations (Table \ref{tab:params_levy}), where the measured exponents are systematically lower than their theoretical values due to finite-time effects in a memoryless process. The vortex system, governed by a quenched velocity field, exhibits persistent motion that inflates the finite-time MSD exponent relative to the expected asymptotic value.

\subsection{Velocity statistics and Holtsmark distribution}

The probability distribution functions (PDFs) of tracer speeds, normalized by their standard deviation $v_c(N)$, show a partial data collapse, suggesting a universal statistical structure dictated by the Biot-Savart law. The tails of these distributions are of particular interest, as Holtsmark theory predicts a power-law decay $P(v) \sim v^{-5/2}$ for the velocity components. Our empirical speed PDFs exhibit heavy tails consistent with this power-law behavior. We observe a slight steepening of the tail as $N$ increases, which may be due to the superposition of multiple velocity fields softening the most extreme velocity events that arise from close encounters with a single filament. Despite this, the overall tail behavior provides strong evidence that the underlying velocity statistics are governed by a L\'evy-stable distribution, as predicted.

\subsection{Velocity correlations and persistence}

To quantify the temporal memory inherent in the deterministic flow, we computed the velocity autocorrelation function (VACF). For the vortex filament system, the VACF exhibits a characteristic decay followed by a region of anti-correlation (a negative lobe), a signature of tracers being deflected by the curved flow fields around filaments. The first zero-crossing time, $\tau_c$, provides a measure of the velocity correlation time. As shown in Table \ref{tab:params_vortex}, $\tau_c$ is approximately 10 s for $N=10$ and $N=40$. For the sparser $N=5$ and $N=20$ cases, the VACF did not cross zero within the simulation window, indicating extremely long correlation times.

The persistence length $L_p$, which measures the typical distance traveled before velocity decorrelation, is found to be 0.313 m for $N=10$ and 0.578 m for $N=40$. These lengths are substantially smaller than the mean inter-filament spacing $\ell_{\text{inter}}$ (9.28 m and 5.85 m, respectively). This implies that velocity memory is lost due to the local curvature of the flow field around a single filament, rather than encounters with distant filaments. This behavior is fundamentally different from a canonical L\'evy walk, which is a memoryless renewal process and thus lacks the anti-correlation signature seen in our VACF analysis.

\subsection{Trapping dynamics and flow topology}

We investigated the relationship between transport dynamics and the local geometry of the flow field using the Okubo-Weiss (OW) parameter. For the $N=40$ case, we found a significant negative correlation (Pearson $r = -0.508$) between the tracer speed and the OW parameter (Table \ref{tab:params_ow}). Since $OW < 0$ in rotation-dominated regions and $OW > 0$ in strain-dominated regions, this result demonstrates that low-speed trapping events occur preferentially in regions of high vorticity, such as the cores of the flow structures induced by the filaments. Conversely, high-speed flights are associated with strain-dominated regions between filaments. A complementary analysis revealed a positive correlation ($r = +0.425$) between speed and the local variance of speed, indicating that the fastest parts of the trajectory are also the most dynamically variable, consistent with rapid motion through the strong velocity gradients near filaments.

The statistics of these low-speed trapping events were characterized by the residence time distribution, $P(\tau)$, shown in Figure \ref{fig:image0}. The distributions for all filament densities exhibit heavy, power-law-like tails. This is qualitatively consistent with the theoretical prediction for a system with Holtsmark velocity statistics, which implies a residence time distribution of $P(\tau) \sim \tau^{-5/2}$ (equivalent to $\tau^{-2.5}$). While the limited statistics of our simulations preclude a precise quantitative fit of the tail exponent, the observed behavior strongly supports the picture of anomalous transport arising from long trapping events in rotation-dominated flow regions.

\begin{figure}[h!]
    \centering
    \includegraphics[width=\columnwidth]{../input_files/plots/step_6_plot_5.png}
    \caption{Residence time distributions, P($\tau$), for low-speed trapping events, defined as speeds below the 25th percentile of the speed distribution, for different numbers of filaments (N). The empirical distributions exhibit heavy-tailed power-law characteristics, qualitatively consistent with the theoretical Holtsmark prediction, P($\tau$) $\sim$ $\tau^{-2.5}$, which is shown as a dashed reference line.}
    \label{fig:image0}
\end{figure}

\subsection{Displacement distributions and non-Gaussianity}

The nature of the transport is further elucidated by the probability distribution of tracer displacements, $P(\Delta x)$. For the $N=40$ configuration, we observed a clear evolution in the shape of this distribution with increasing lag time $\Delta t$. At short lag times ($\Delta t = 1$ s), the distribution is highly non-Gaussian, with pronounced heavy tails that are well-described by a L\'evy-stable distribution with a stability index consistent with the Holtsmark prediction ($\alpha_{\text{stable}} \approx 1.5$). As the lag time increases towards and beyond the velocity correlation time ($\tau_c \approx 9.5$ s), the central part of the distribution progressively approaches a Gaussian form, as expected from the central limit theorem acting on increasingly decorrelated steps. However, the tails remain heavy, indicating that the transport process retains its L\'evy-like character over all observed timescales.

\subsection{Signatures of non-ergodicity}

The quenched nature of the velocity field leads to signatures of non-ergodic behavior. We observed a pronounced spread in the TAMSD curves calculated for individual tracers within the same filament configuration. This realization-to-realization variability is not statistical noise but a physical consequence of the fact that each tracer samples a different, fixed spatial region of the flow field. The spread is largest for the sparse $N=5$ case and narrows as $N$ increases, as the velocity field becomes more spatially homogeneous. However, even at $N=40$, significant spread persists, indicating that the observation time is insufficient for a single trajectory to fully sample the statistical properties of the entire disordered environment. This weak ergodicity breaking is a hallmark of transport in quenched media with long-range correlations and heavy-tailed statistics.

\end{document}
                